\section{Research baseline and contribution map}
\label{sec:introduction}

The manuscript is a developing research corpus, with newer incoming
reports that sometimes settle questions still listed as open elsewhere.
A continuation must reconcile that history before selecting targets.
We therefore use a single pinned version of Vladimir Reshetnikov's
ProveIt repository \cite{repo}: commit \snapshot, dated 10 October 2026
UTC. The inspection covered the manuscript's 40 chapter source files,
its editorial and validation material, and all 16 archives present in
\path{docs/incoming}. The 65 retrieved source objects, including those
archives, have paths, byte counts, and SHA-256 hashes in the accompanying
provenance manifest. This is a source and dependency audit of the present
continuation, not an assertion that every proof in the entire corpus
has been reverified.

The most useful unresolved targets lie at the intersection of positive
integral representations, exact formal relations, and zero geometry.
The work below develops all three. Its principal resolved conjecture is
the sharp critical constant in the real-order report \cite{realorder}.
The prime-order distribution calculation answers a proposed extension
of reflection torsion \cite{integraljets}. The Lerch results extend the
branch analysis beyond the cases already proved in the incoming
boundary and global-phase reports \cite{lershape:Boundary,lershape:Globalphase}.
The $S_6$ computation gives a rigorous limitation of a larger finite
proof search, while leaving the numerical identity open.

\subsection{Conventions and distinctions}

We use decreasing-index multiple polylogarithms:
\[
 \Li_{s_1,\ldots,s_d}(z_1,\ldots,z_d)
 =\sum_{n_1>\cdots>n_d>0}
   \frac{z_1^{n_1}\cdots z_d^{n_d}}
        {n_1^{s_1}\cdots n_d^{s_d}},
\]
initially in a domain of convergence. In particular
$F_{a,b}(z)=\Li_{a,b}(z,1)$ and $g_{a,b}=\Im F_{a,b}(i)$.
The positive real orders $a,b$ need not be integers. At boundary points
we specify analytic or Abel values when the ordinary series diverges.
We write $\beta(s)=\sum_{j\ge0}(-1)^j(2j+1)^{-s}$ and
$\eta(s)=\sum_{j\ge1}(-1)^{j-1}j^{-s}$ where their series converge,
and use their analytic continuations elsewhere. Standard conventions
are also recorded in the DLMF \cite{dlmf-polylog,dlmf-hurwitz}.

There are two different weights in this article: $a+b$ is the total
order of a Gaussian double polylogarithm, whereas the integer coefficients
$a_p$ in a distribution relation are arbitrary formal parameters.
The parameters $n,k,\rho$ in the Lerch family denote Stieltjes index,
derivative order, and deformation parameter. They are introduced afresh
in Section~\ref{lershape:sec:main}.

An equality of complex periods, membership in a specified rational
relation span, and a Smith invariant of an integral formal module are
different assertions. Our notation and proofs retain those distinctions.
In particular, the torsion theorem does not claim torsion among complex
numbers, and the $S_6$ separator does not claim numerical independence.

\subsection{What is added to the inspected baseline}

\begin{table}[htbp]
\centering
\begin{tabular}{@{}L{0.21\textwidth}L{0.39\textwidth}L{0.33\textwidth}@{}}
\toprule
Topic & Result in this continuation & Scope and proof status\\
\midrule
Critical Gaussian values & Strict allocation monotonicity for every
positive total order; sharp Euler constant on $a+b=1$
& Analytic proofs; resolves the incoming critical conjecture.\\[5pt]
Global envelope & A unique maximum of $\beta(w)+2^{-w}\eta(w)$,
sharp bound for all positive Gaussian orders, and inverse coefficients
& Analytic shape theorem plus rational maximum enclosure.\\[5pt]
Euler remainders & A positive secant integral and an exact example of
increasing normalized errors off the critical line
& Analytic identities; prevents an unsupported extension in $N$.\\[5pt]
Cyclic distributions & Integral coinvariants and all one-variable
finite jets for a split prime-order action
& Complete formula under explicit arithmetic hypotheses.\\[5pt]
Lerch zeros & Finite extremal-zero comparison; lower-branch dichotomy;
order-two minimum and order-three decrease
& Analytic theorems with exact endpoint sign certificates.\\[5pt]
Weight-seven reduction & Separation of the $S_6$ target from 5,131
specified relations through depth four
& Integer certificate; the period identity remains conjectural.\\
\bottomrule
\end{tabular}
\end{table}

The beta allocation identity is the main analytic device. It isolates
the split between the two orders in a beta probability law, leaving a
strictly decreasing scalar kernel. A covariance then proves
monotonicity. Optimizing the endpoint kernel reduces a two-parameter
extremum to a one-variable Dirichlet-function problem. Its strict
log-concavity is an application of a classical normalized Mellin
principle \cite{flm}, for which we supply the strict proof actually needed.

The distribution calculation follows a different route: a finite
integral resolution removes moving orbits from a periodic cohomology
complex, and the remaining fixed part becomes a Koszul calculation.
The resulting character condition explains how weights on different
primes can constrain one another. This is more informative than testing
each weight separately. Classical universal-distribution work is an
essential antecedent, explicitly credited in the theorem section.

For Lerch zeros, positivity acts after a resolvent comparison rather
than directly on the original sign-changing kernel. Beyond the largest
zero, every shifted summand has a known sign, so a finite comparison
works even at parameters where some zeros collide. At index two,
the three-block sign pattern of the coefficient sequence additionally
restricts all possible stationary points of the lower zero branch.

\subsection{Results deliberately retained as prior work}

The positive double integral, the sharp finite signed-measure threshold
$a+b\ge1$, and the endpoint atom on $a+b=1$ are already present in
\cite{realorder}. The current manuscript already proves the $S_4$
identity; it is not a conjecture to rediscover. Integral freeness,
scalar reflection torsion, and reflection jets are already proved in
\cite{integraljets}. The index-two Lerch zero count and the unique
minimum at derivative order one are also prior results
\cite{lershape:Boundary}. We recall short arguments where they support
new theorems, with explicit attribution.

The references are chosen to support the actual proof mechanisms and
to identify nearby literature. They do not establish that every result
here is new in the full mathematical literature. The specific
continuation claims concern the inspected repository baseline.
